\chapter{系统实现与训练优化}

\section{开发环境与依赖配置}

系统开发基于以下环境：

（1）操作系统：Windows 11 或 Linux

（2）Python 版本：3.x

（3）深度学习框架：PyTorch 2.x

（4）数据处理库：NumPy

（5）可视化工具：Matplotlib

\section{数据加载模块实现}

StaticPreprocessor 类完成数据预处理的全部流程，具体包括：

（1）读取并解析原始 CSV 文件。

（2）清洗数据，剔除异常记录。

（3）构建滑动时间窗口。

（4）对特征进行标准化处理。

（5）将处理后的数据保存为 NPZ 格式，主要字段包括 X\_train、X\_train\_spatial、y\_train 等。

数据划分由两个类实现：IndomainSplitter 负责内域划分，OutdomainSplitter 负责外域划分。

\section{模型实现细节}

STCAModel 类继承自 PyTorch 的 nn.Module，核心是前向传播逻辑，包含五个处理阶段：

（1）空间编码器接收空间输入，生成空间嵌入向量。

（2）时序编码器采用 LSTM 处理时序输入，输出时序上下文向量。

（3）交叉注意力模块将时空特征进行融合。

（4）稀疏表示层从融合特征中提取紧凑表示。

（5）分类器最终输出 NLOS 信号的识别概率。

模型参数采用 Xavier 均匀分布进行初始化。

\section{训练器实现}

训练器将训练循环和评估逻辑封装在两个核心方法中。

fit() 方法负责训练过程：

（1）每个 epoch 遍历训练集中的所有批次。

（2）执行前向传播，计算损失值。

（3）执行反向传播，更新模型参数。

（4）根据调度策略调整学习率。

（5）在验证集上评估模型，保存验证准确率最高时的权重。

evaluate() 方法负责模型评估：

（1）遍历测试集中的所有批次。

（2）执行前向传播，获取预测结果。

（3）计算准确率、精确率、召回率等性能指标。

\section{训练优化策略}

训练过程中采用以下优化措施：

（1）使用 Adam 优化器，设置 $\beta_1=0.9$、$\beta_2=0.999$。

（2）学习率初始值为$10^{-4}$，采用衰减因子 0.5 的动态调度策略。

（3）梯度裁剪阈值设为 1.0，避免梯度爆炸。

（4）当验证集指标持续不提升时，触发早停机制终止训练。

\section{超参数确定实验}

为确定最优模型超参数配置，本研究进行了系统的参数搜索实验。表~\ref{tab:parameter_search} 展示了不同参数配置下的模型性能对比结果。

\begin{table}[h]
\centering
\caption{超参数配置搜索实验结果}
\label{tab:parameter_search}
\begin{tabular}{cccccc}
\hline
空间编码器 & 时序编码器 & 交叉注意力 & 分类器 & Loss & 准确率 \\
\hline
(16,1,1,32,0.5) & (16,1,0.5) & (16,1,0.5) & (16,8,0.5) & 0.7846 & 69.96\% \\
(16,2,1,32,0.5) & (16,1,0.5) & (16,1,0.5) & (16,8,0.5) & 0.7959 & 70.11\% \\
(16,4,1,32,0.5) & (16,1,0.5) & (16,1,0.5) & (16,8,0.5) & 0.8193 & 70.22\% \\
(8,1,1,16,0.5) & (8,1,0.5) & (8,1,0.5) & (8,4,0.5) & 0.6273 & 73.89\% \\
(32,1,1,64,0.5) & (32,1,0.5) & (32,1,0.5) & (32,16,0.5) & 1.3517 & 60.97\% \\
(16,1,1,32,0.5) & (16,2,0.5) & (16,1,0.5) & (16,8,0.5) & 0.8471 & 69.39\% \\
(16,1,1,32,0.5) & (16,1,0.5) & (16,1,0.5) & (16,8,0.5) & 0.8169 & 68.88\% \\
(16,1,2,32,0.5) & (16,1,0.5) & (16,1,0.5) & (16,8,0.5) & 0.8233 & 69.35\% \\
\hline
\end{tabular}
\end{table}

表注：空间编码器参数格式为 (embed\_dim, num\_layers, num\_heads, d\_ff, dropout)；时序编码器为 (embed\_dim, num\_layers, dropout)；交叉注意力为 (embed\_dim, num\_heads, dropout)；分类器为 (hidden1, hidden2, dropout)。

实验结果表明，当空间编码器采用 1 层、1 注意力头、前馈网络维度 32、Dropout 率 0.5，时序编码器采用 1 层 LSTM、Dropout 率 0.5，交叉注意力模块采用 1 注意力头、Dropout 率 0.5，分类器隐藏层维度为 [16, 8]、Dropout 率 0.5 时，模型获得最优性能，Loss 为 0.6273，准确率为 73.89\%。该配置下各模块容量适配，既避免了欠拟合，也有效抑制了过拟合。

\section{模块实现细节}

\subsection{空间编码器模块}

空间编码器在 SpatialEncoder 类中实现，位于 modules/spatial\_encoder.py 文件。核心代码结构如下：

\begin{verbatim}
class SpatialEncoder(nn.Module):
    def __init__(self, input_dim, embed_dim, num_heads, 
                 num_layers, d_ff, dropout_rate):
        super().__init__()
        self.feature_embed = nn.Sequential(
            nn.Linear(input_dim, embed_dim),
            nn.ReLU(),
            nn.Dropout(dropout_rate)
        )
        self.attention = nn.MultiheadAttention(
            embed_dim, num_heads, dropout=dropout_rate
        )
\end{verbatim}

空间编码器的核心是自注意力机制，通过计算卫星之间的相关性权重，自动学习哪些卫星的特征更重要。对于变长卫星输入，采用掩码（mask）机制处理填充位置。

\subsection{时序编码器模块}

时序编码器在 TemporalEncoder 类中实现，位于 modules/temporal\_encoder.py 文件。采用 PyTorch 内置的 LSTM 模块：

\begin{verbatim}
class TemporalEncoder(nn.Module):
    def __init__(self, input_dim, embed_dim, num_layers, 
                 dropout_rate, bidirectional):
        super().__init__()
        self.lstm = nn.LSTM(
            input_size=input_dim,
            hidden_size=embed_dim,
            num_layers=num_layers,
            batch_first=True,
            bidirectional=bidirectional,
            dropout=dropout_rate if num_layers > 1 else 0
        )
\end{verbatim}

双向 LSTM 的输出维度为 embed\_dim $\times$ 2，本研究中为 32 维。提取最后一个时间步的隐藏状态作为时序上下文向量。

\subsection{交叉注意力模块}

交叉注意力模块在 CrossAttention 类中实现，位于 modules/cross\_attention.py 文件：

\begin{verbatim}
class CrossAttention(nn.Module):
    def __init__(self, embed_dim, num_heads, dropout_rate,
                 query_input_dim, kv_input_dim):
        super().__init__()
        self.query_proj = nn.Linear(query_input_dim, embed_dim)
        self.key_proj = nn.Linear(kv_input_dim, embed_dim)
        self.value_proj = nn.Linear(kv_input_dim, embed_dim)
        self.attention = nn.MultiheadAttention(
            embed_dim, num_heads, dropout=dropout_rate
        )
\end{verbatim}

交叉注意力的 Query 来自时序编码器输出，Key 和 Value 来自空间编码器输出。这种设计使模型能够根据时序状态，自适应地关注相关的空间特征。

\section{训练流程图}

图~\ref{fig:training-flowchart} 展示了 STCA 模型的完整训练流程。

\begin{figure}[h]
\centering
\fbox{
\begin{minipage}{0.9\linewidth}
\centering
\begin{enumerate}
    \item 加载 NPZ 格式数据
    \item 创建 DataLoader，设置 batch 大小
    \item 初始化模型、优化器、损失函数
    \item \textbf{For} epoch = 1 \textbf{to} num\_epochs \textbf{do}
    \begin{itemize}
        \item \textbf{For} batch \textbf{in} DataLoader \textbf{do}
        \begin{itemize}
            \item 前向传播：计算预测值
            \item 计算损失：BCELoss
            \item 反向传播：更新参数
        \end{itemize}
        \item 验证集评估
        \item 学习率调整
        \item 早停检查
    \end{itemize}
    \item 保存最佳模型权重
\end{enumerate}
\end{minipage}
}
\caption{STCA 模型训练流程图}
\label{fig:training-flowchart}
\end{figure}

\section{超参数调优过程}

超参数调优采用网格搜索与经验法则相结合的策略：

（1）嵌入维度：参考 Zeng 等人（2024）的论文推荐值，从 16 开始搜索，测试 16、32、64 三种配置。实验表明，16 维嵌入在 GNSS NLOS 识别任务上已获得较好性能，继续增加维度带来的性能提升有限，反而增加计算负担。

（2）层数：空间编码器测试 1、2、4 层，时序编码器测试 1、2 层。过深的网络容易导致过拟合，且增加训练时间。

（3）Dropout 率：测试 0.3、0.5、0.6 三种配置。0.5 在抑制过拟合和保留信息之间取得最佳平衡。

（4）学习率：测试$10^{-3}$、$10^{-4}$、$10^{-5}$三种配置。$10^{-4}$收敛速度适中，最终性能最优。

（5）批次大小：测试 8、16、32 三种配置。16 在 GPU 内存利用率和梯度稳定性之间取得平衡。

\section{本章小结}

本章介绍了 STCA 模型的系统实现细节。开发环境采用 PyTorch 2.x 深度学习框架。数据加载模块实现数据预处理的全部流程。模型实现包含空间编码器、时序编码器、交叉注意力模块、稀疏表示层和分类器五个模块。训练优化策略采用 Adam 优化器、动态学习率调度和早停机制。超参数确定实验表明，空间编码器 1 层、时序编码器 1 层、嵌入维度 16/32 的配置最优。
